% Overall technical route for Section 3.2, rebuilt 2026-10-29 from the
% confirmed v2 preview.  Three reconstruction studies co-build one shared
% representation; the fourth study reads that representation at t0, propagates
% a given action through typed relations and integrates a learned dynamics.
% Subsequent recorded observations are written back, and the write-back
% bypasses the initialisation step.
\begingroup
% Shared editable vector language for the five Section 3.2 technical routes.
%
% The figures are explanatory diagrams, not experimental results.  A single
% visual grammar is used throughout:
%   * a dark header band names a block; white type is reserved for headers;
%   * a numbered marker gives the reading order of a route, and the marker of a
%     stage that carries the study's own mechanism is filled;
%   * white cards with a thin theme-coloured rule hold stage titles and bodies;
%   * tinted panels hold mechanism sketches, and a heavier frame marks the
%     representation or relation that is handed to the next study;
%   * solid arrows are information flow, dashed arrows are observation
%     feedback or a model-predicted interval.
% Colour is always paired with a marker, line style or label so that the
% grayscale proof remains readable.
\linespread{1}\selectfont
\usetikzlibrary{arrows.meta,calc,positioning,matrix,decorations.pathreplacing}

% ---- palette -------------------------------------------------------------
\definecolor{trBlueDark}{HTML}{1F4E79}
\definecolor{trBlueMid}{HTML}{2F6690}
\definecolor{trBlueSoft}{HTML}{7FA6C4}
\definecolor{trBluePale}{HTML}{E8F1F8}
\definecolor{trBlueWash}{HTML}{F4F8FB}
\definecolor{trGreenDark}{HTML}{2E5E40}
\definecolor{trGreenMid}{HTML}{3E7A52}
\definecolor{trGreenSoft}{HTML}{8FB39A}
\definecolor{trGreenPale}{HTML}{E9F2EB}
\definecolor{trGreenWash}{HTML}{F4F9F5}
\definecolor{trOchreDark}{HTML}{7A5018}
\definecolor{trOchreMid}{HTML}{9E6B27}
\definecolor{trOchreSoft}{HTML}{C9A46B}
\definecolor{trOchrePale}{HTML}{F8EEDC}
\definecolor{trOchreWash}{HTML}{FCF7EE}
\definecolor{trVioletDark}{HTML}{4A3A6A}
\definecolor{trVioletMid}{HTML}{67528E}
\definecolor{trVioletSoft}{HTML}{A294BE}
\definecolor{trVioletPale}{HTML}{EFEAF6}
\definecolor{trVioletWash}{HTML}{F8F5FC}
\definecolor{trInkSoft}{HTML}{4D4D4D}
\definecolor{trRule}{HTML}{9AA5AE}

\newcommand{\trTheme}[1]{%
  \colorlet{trDark}{tr#1Dark}%
  \colorlet{trMid}{tr#1Mid}%
  \colorlet{trSoft}{tr#1Soft}%
  \colorlet{trPale}{tr#1Pale}%
  \colorlet{trFill}{tr#1Pale}%
  \colorlet{trWash}{tr#1Wash}%
}
\trTheme{Blue}

% ---- shared type scale for every Part 3 figure ----------------------------
% Set sizes here rather than overriding individual diagrams.  The report uses
% the natural figure width, so these remain the printed sizes; subscripts use
% the normal mathematical script sizes.
\newcommand{\trFontHead}{\songti\bfseries\fontsize{9.2pt}{11.2pt}\selectfont}
\newcommand{\trFontSubhead}{\songti\bfseries\fontsize{8.2pt}{10.2pt}\selectfont}
\newcommand{\trFontBody}{\songti\fontsize{8.2pt}{10.2pt}\selectfont}
\newcommand{\trFontSmall}{\songti\fontsize{8pt}{10pt}\selectfont}
\newcommand{\trFontTiny}{\songti\fontsize{8pt}{10pt}\selectfont}
\newcommand{\trFontMath}{\fontsize{9.2pt}{11.2pt}\selectfont}
\newcommand{\trFontMathSmall}{\fontsize{8.5pt}{10.4pt}\selectfont}
\newcommand{\trFontMarker}{\songti\bfseries\fontsize{8pt}{10pt}\selectfont}
\newcommand{\trFontSubcaption}{\songti\mdseries\fontsize{9.2pt}{11.2pt}\selectfont}

% ---- base styles ---------------------------------------------------------
\tikzset{
  trBase/.style={x=1cm,y=1cm,text=black,line cap=round,line join=round,
    font=\trFontBody,
    every node/.style={align=center,outer sep=0pt,inner sep=0pt}},
  trTitle/.style={font=\trFontHead},
  trTitleSmall/.style={font=\trFontSubhead},
  trHead/.style={font=\trFontHead,text=white},
  trBody/.style={font=\trFontBody},
  trNote/.style={font=\trFontSmall},
  trTiny/.style={font=\trFontTiny},
  trLabel/.style={font=\trFontSmall},
  trMath/.style={font=\trFontMath},
  trMathSmall/.style={font=\trFontMathSmall},
  trNum/.style={circle,draw=trDark,fill=trDark,text=white,
    line width=0.5pt,inner sep=0pt,minimum size=0.36cm,
    font=\trFontMarker},
  trNumOpen/.style={trNum,fill=white,text=trDark},
  trCard/.style={fill=white,draw=trMid,line width=0.55pt,rounded corners=1.2pt},
  trFrame/.style={draw=black,line width=0.65pt},
  trArrow/.style={draw=trDark,line width=0.8pt,
    -{Latex[length=1.7mm,width=1.15mm]}},
  trArrowThin/.style={draw=trDark,line width=0.6pt,
    -{Latex[length=1.4mm,width=0.95mm]}},
  trFlow/.style={draw=black!72,line width=0.7pt,
    -{Latex[length=1.6mm,width=1.05mm]}},
  trFeedback/.style={draw=trDark,line width=0.8pt,dash pattern=on 3pt off 2pt,
    -{Latex[length=1.7mm,width=1.15mm]}},
  trPredict/.style={draw=trMid,line width=0.75pt,dash pattern=on 2.2pt off 1.6pt,
    -{Latex[length=1.6mm,width=1.05mm]}},
  trGuide/.style={draw=black!35,line width=0.45pt},
  trTag/.style={draw=trMid,fill=white,line width=0.45pt,inner sep=2.2pt,
    rounded corners=0.8pt,font=\trFontSmall},
  trAccent/.style={draw=trMid,line width=0.9pt},
  trAccentThin/.style={draw=trMid,line width=0.6pt},
  trPossible/.style={draw=trSoft,line width=0.7pt,dash pattern=on 1.8pt off 1.4pt},
  trTimeLine/.style={draw=black!55,line width=0.6pt,
    -{Latex[length=1.5mm,width=0.95mm]}},
  trDown/.style={draw=trDark,line width=0.8pt,
    -{Latex[length=1.6mm,width=1.05mm]}},
  trDouble/.style={draw=trDark,line width=0.8pt,
    {Latex[length=1.5mm,width=1.0mm]}-{Latex[length=1.5mm,width=1.0mm]}},
  trDoubleSmall/.style={draw=trDark,line width=0.7pt,
    {Latex[length=1.1mm,width=0.75mm]}-{Latex[length=1.1mm,width=0.75mm]}}
}

% ---- reusable pieces -----------------------------------------------------
% Header band.  #1 left, #2 right, #3 bottom, #4 top, #5 text.
\newcommand{\trHeader}[5]{%
  \path[fill=trDark] (#1,#3) rectangle (#2,#4);
  \node[trHead] at ({#1/2+#2/2},{#3/2+#4/2}) {#5}
}

% ---- compact detailed-route primitives (2026-10-30) ----------------------
% The four detailed routes use a downward y axis so that a route reads from
% the input strip at the top to the output strip at the bottom.  Their canvases
% are 15.56 cm wide and only as tall as the content; the styles below keep the
% four redraws on the same type scale as the overview (Fig. 11).
\tikzset{
  trRouteBase/.style={x=1cm,y=-1cm,text=black,line cap=round,line join=round,
    font=\trFontBody,
    every node/.style={align=center,outer sep=0pt,inner sep=0pt}},
  trRouteHead/.style={font=\trFontHead,text=white},
  trRouteBody/.style={font=\trFontBody},
  trRouteSmall/.style={font=\trFontSmall},
  trRouteTiny/.style={font=\trFontTiny},
  trRouteMath/.style={font=\trFontMath},
  trRouteMathSmall/.style={font=\trFontMathSmall},
  trRouteFlow/.style={draw=trDark,line width=0.72pt,
    -{Latex[length=1.5mm,width=1.0mm]}},
  trRouteThin/.style={draw=trDark,line width=0.6pt,
    -{Latex[length=1.3mm,width=0.9mm]}},
  trRouteBack/.style={draw=trDark,line width=0.65pt,
    dash pattern=on 2.6pt off 1.8pt,-{Latex[length=1.4mm,width=0.95mm]}},
  trRouteRule/.style={draw=trMid,line width=0.55pt},
  trRouteMotion/.style={draw=trBlueDark,line width=0.95pt},
  trRoutePossible/.style={draw=black!52,line width=0.65pt,
    dash pattern=on 1.8pt off 1.4pt},
  trRouteCard/.style={fill=white,draw=trMid,line width=0.6pt},
  trRouteStrip/.style={fill=trWash,draw=trMid,line width=0.55pt},
  trRouteFocus/.style={fill=trPale,draw=trDark,line width=0.7pt}
}

\trTheme{Blue}
\newcommand{\ovCard}[6]{% x1 x2 y1 y2 title subtitle
  \pgfmathsetlengthmacro{\ovCardTextW}{(#2-#1-0.26)*1.06575cm}%
  \path[fill=white,draw=trMid,line width=0.5pt,rounded corners=1pt]
    (#1,#3) rectangle (#2,#4);
  \node[font=\trFontSubhead,
    text width=\ovCardTextW,align=center]
    at ({#1/2+#2/2},{#3/2+#4/2+0.19}) {#5};
  \node[font=\trFontTiny,text=black!78,
    text width=\ovCardTextW,align=center]
    at ({#1/2+#2/2},{#3/2+#4/2-0.18}) {#6};
}
\newcommand{\ovCardOne}[5]{% x1 x2 y1 y2 title
  \pgfmathsetlengthmacro{\ovOneTextW}{(#2-#1-0.24)*1.06575cm}%
  \path[fill=white,draw=trMid,line width=0.5pt,rounded corners=1pt]
    (#1,#3) rectangle (#2,#4);
  \node[font=\trFontSubhead,
    text width=\ovOneTextW,align=center]
    at ({#1/2+#2/2},{#3/2+#4/2}) {#5};
}
\newcommand{\ovActionCard}[6]{% x1 x2 y1 y2 title subtitle
  \pgfmathsetlengthmacro{\ovActionTextW}{(#2-#1-0.18)*1.06575cm}%
  \path[fill=white,draw=trMid,line width=0.5pt,rounded corners=1pt]
    (#1,#3) rectangle (#2,#4);
  \node[font=\trFontSubhead,
    text width=\ovActionTextW,align=center]
    at ({#1/2+#2/2},{#4-0.25}) {#5};
  \node[font=\trFontTiny,text=black!78,
    text width=\ovActionTextW,align=center]
    at ({#1/2+#2/2},{#3+0.25}) {#6};
}
\newcommand{\ovActionCardOne}[5]{% x1 x2 y1 y2 title
  \pgfmathsetlengthmacro{\ovActionOneTextW}{(#2-#1-0.18)*1.06575cm}%
  \path[fill=white,draw=trMid,line width=0.5pt,rounded corners=1pt]
    (#1,#3) rectangle (#2,#4);
  \node[font=\trFontSubhead,
    text width=\ovActionOneTextW,align=center]
    at ({#1/2+#2/2},{#3/2+#4/2}) {#5};
}
\begin{tikzpicture}[trBase,x=1.06575cm,yscale=1.14]
  \path[use as bounding box] (0.00,-1.28) rectangle (14.60,12.46);

  % ================= video observation =================
  \path[fill=trBluePale,draw=trBlueMid,line width=0.7pt,rounded corners=1.4pt]
    (0.25,11.55) rectangle (13.55,12.41);
  \node[font=\trFontHead] at (6.90,12.11)
    {视频观测与可用几何线索};
  \node[font=\trFontSmall,text=black!80]
    at (6.90,11.74) {视频帧 · 相机信息 · 跨时刻对应 · 可用深度};
  \draw[trDown] (6.90,11.52) -- (6.90,11.08);

  % ================= collaborative construction =================
  \path[fill=trBlueWash,draw=trBlueDark,line width=1.0pt]
    (0.25,3.05) rectangle (13.55,11.05);
  \trHeader{0.25}{13.55}{10.60}{11.05}{连续动态表征的协同构建};

  % Three study columns.  The header band carries the study number and name;
  % the three cards below are the route stages; the lower panel is a mechanism
  % sketch, not a decorative picture.
  % ---- column 1 ----
  \path[fill=white,draw=trBlueMid!55,line width=0.55pt]
    (0.42,4.35) rectangle (4.52,10.50);
  \path[fill=trBlueDark] (0.42,9.58) rectangle (4.52,10.50);
  \node[trNumOpen,fill=white,text=trBlueDark,minimum size=0.46cm] at (0.83,10.04) {1};
  \node[font=\trFontHead,text=white]
    at (2.78,10.04) {不确定性感知的\\连续时间场景表征};
  \ovCard{0.55}{4.39}{8.70}{9.50}{观测置信度估计}{循环一致性、重投影与可见性}
  \ovCard{0.55}{4.39}{7.84}{8.64}{时间支撑与表达配置}{局部有效区间与自由度}
  \ovCard{0.55}{4.39}{6.98}{7.78}{几何—运动联合求解}{轨迹估计与几何回校正}

  % sketch: discrete evidence and a continuous trajectory
  \path[fill=trBluePale!55,draw=trBlueMid!45,line width=0.5pt]
    (0.55,4.48) rectangle (4.39,6.92);
  \node[trTitleSmall,anchor=north west] at (0.68,6.82) {离散证据与连续轨迹};
  \node[trLabel,anchor=west] at (0.68,4.64) {$\bullet$ 高置信度\quad$\circ$ 低置信度};
  % support band and trajectory
  \draw[trBlueMid,line width=8pt,opacity=0.15,line cap=round]
    plot[smooth] coordinates {(0.78,5.62) (1.32,5.90) (1.88,5.96) (2.42,5.70)
      (2.96,5.94) (3.50,6.20) (4.10,6.06)};
  \draw[trBlueDark,line width=1.1pt]
    plot[smooth] coordinates {(0.78,5.62) (1.32,5.90) (1.88,5.96) (2.42,5.70)
      (2.96,5.94) (3.50,6.20) (4.10,6.06)};
  \foreach \xx/\yy in {0.96/5.72,1.66/5.94,2.46/5.73,3.26/6.03,4.02/6.08}{
    \fill[black] (\xx,\yy) circle (0.055);
  }
  \draw[black,line width=0.7pt,fill=white] (2.06,5.86) circle (0.062);
  \draw[-{Latex[length=1.6mm,width=1.0mm]},draw=trVioletMid,line width=0.9pt]
    (2.72,5.92) -- (3.10,6.34);
  \node[font=\trFontMathSmall,text=trVioletMid]
    at (3.24,6.39) {$\mathbf v(t)$};
  \draw[black!70,line width=0.5pt] (1.55,5.30) -- (1.55,5.22)
    (1.55,5.22) -- (3.32,5.22) (3.32,5.22) -- (3.32,5.30);
  \node[trLabel] at (2.44,5.06) {局部时间支撑};

  % ---- column 2 ----
  \path[fill=white,draw=trGreenMid!55,line width=0.55pt]
    (4.80,4.35) rectangle (8.90,10.50);
  \path[fill=trGreenDark] (4.80,9.58) rectangle (8.90,10.50);
  \node[trNumOpen,fill=white,text=trGreenDark,minimum size=0.46cm] at (5.21,10.04) {2};
  \node[font=\trFontHead,text=white]
    at (7.16,10.04) {连续运动引导的\\空间结构表征};
  \ovCard{4.93}{8.77}{8.70}{9.50}{多证据关联}{运动、几何、外观与可见性}
  \ovCard{4.93}{8.77}{7.84}{8.64}{部件软归属与共享门控}{归属、待定关系与共享范围}
  \ovCard{4.93}{8.77}{6.98}{7.78}{归属—运动联合恢复}{共享范围与跨时刻关联}

  % sketch: two parts, motion vectors and an undecided unit
  \path[fill=trGreenPale!55,draw=trGreenMid!45,line width=0.5pt]
    (4.93,4.48) rectangle (8.77,6.92);
  \node[trTitleSmall,anchor=north west] at (5.06,6.82) {归属约束共享范围};
  \foreach \cx/\cy/\col in {5.92/5.68/trBlueMid,7.72/5.80/trGreenMid}{
    \draw[\col,dash pattern=on 2pt off 1.4pt,line width=0.6pt,fill=white,opacity=0.85]
      (\cx,\cy) ellipse (0.74 and 0.46);
  }
  \foreach \xx/\yy in {5.58/5.54,5.94/5.84,6.28/5.60,5.84/5.38}{
    \draw[trBlueDark,line width=0.75pt,-{Latex[length=1.2mm,width=0.8mm]}]
      (\xx,\yy) -- ++(0.14,0.17);
    \fill[trBlueDark] (\xx,\yy) circle (0.048);
  }
  \foreach \xx/\yy in {7.38/5.70,7.76/5.94,8.10/5.72,7.68/5.50}{
    \draw[trGreenDark,line width=0.75pt,-{Latex[length=1.2mm,width=0.8mm]}]
      (\xx,\yy) -- ++(0.14,0.17);
    \draw[trGreenDark,line width=0.65pt] (\xx-0.048,\yy-0.048) rectangle ++(0.096,0.096);
  }
  \node[trLabel] at (5.92,5.02) {部件 A};
  \node[trLabel] at (7.72,5.02) {部件 B};
  \draw[trGreenDark,line width=0.65pt,fill=white] (6.83,6.12)
    -- (6.95,6.24) -- (6.83,6.36) -- (6.71,6.24) -- cycle;
  \draw[black!55,line width=0.55pt,dash pattern=on 1.6pt off 1.2pt]
    (6.74,6.10) -- (6.22,5.90) (6.92,6.10) -- (7.44,5.96);
  \node[trLabel,anchor=east] at (8.60,6.48) {待定归属};
  \draw[black!55,line width=0.45pt] (7.48,6.40) -- (6.96,6.28);

  % ---- column 3 ----
  \path[fill=white,draw=trOchreMid!55,line width=0.55pt]
    (9.18,4.35) rectangle (13.28,10.50);
  \path[fill=trOchreDark] (9.18,9.58) rectangle (13.28,10.50);
  \node[trNumOpen,fill=white,text=trOchreDark,minimum size=0.46cm] at (9.59,10.04) {3};
  \node[font=\trFontHead,text=white]
    at (11.54,10.04) {物理连接变化下的\\连续动态场景表征};
  \ovCard{9.31}{13.15}{8.70}{9.50}{物理连接判断}{观测证据与可用事件提示}
  \ovCard{9.31}{13.15}{7.84}{8.64}{局部连接与表面调整}{解除失效联系、更新表面}
  \ovCard{9.31}{13.15}{6.98}{7.78}{材料运动与表面历史维护}{稳定区历史、新生边界}

  % sketch: before/after separation and two histories
  \path[fill=trOchrePale!55,draw=trOchreMid!45,line width=0.5pt]
    (9.31,4.48) rectangle (13.15,6.92);
  \node[font=\trFontSubhead,anchor=north west] at (9.44,6.86) {连接变化与表面历史};
  % intact slab
  \path[fill=trBluePale,draw=trBlueDark,line width=0.6pt]
    (9.58,5.88) -- (9.95,6.10) -- (10.58,6.10) -- (10.58,5.80) --
    (10.21,5.60) -- (9.58,5.60) -- cycle;
  \path[fill=trBlueSoft!35,draw=trBlueDark,line width=0.6pt]
    (9.58,5.60) -- (10.21,5.60) -- (10.21,5.45) -- (9.58,5.45) -- cycle;
  \draw[trBlueDark,line width=0.6pt] (10.21,5.60) -- (10.58,5.80);
  \draw[trOchreMid,line width=0.55pt,dash pattern=on 1.8pt off 1.2pt]
    (10.30,5.60) -- (10.30,6.10);
  % arrow and event
  \draw[trFlow] (10.72,5.84) -- (11.06,5.84);
  \node[font=\trFontTiny,fill=trOchrePale!55,inner sep=0.6pt]
    at (10.90,5.66) {断裂};
  \draw[trOchreDark,line width=0.9pt] (11.36,5.36) -- (11.36,6.28);
  % separated pieces with new inner faces
  \path[fill=trBluePale,draw=trBlueDark,line width=0.6pt]
    (11.52,5.92) -- (11.82,6.10) -- (12.10,5.92) -- (12.10,5.64) --
    (11.82,5.48) -- (11.52,5.64) -- cycle;
  \path[fill=trOchrePale,draw=trOchreDark,line width=0.8pt]
    (12.10,5.92) -- (11.82,6.10) -- (11.82,5.48) -- (12.10,5.64) -- cycle;
  \path[fill=trBluePale,draw=trBlueDark,line width=0.6pt]
    (12.30,5.92) -- (12.58,6.10) -- (12.86,5.92) -- (12.86,5.64) --
    (12.58,5.48) -- (12.30,5.64) -- cycle;
  \path[fill=trOchrePale,draw=trOchreDark,line width=0.8pt]
    (12.30,5.92) -- (12.58,6.10) -- (12.58,5.48) -- (12.30,5.64) -- cycle;
  \node[font=\trFontTiny,text=trOchreDark]
    at (12.55,6.35) {新生边界};
  \draw[trOchreDark,line width=0.35pt] (12.55,6.22) -- (12.58,5.98);
  % Short explanatory lines keep the event sketch readable at report scale.
  \node[font=\trFontSmall,anchor=west] at (9.44,5.18)
    {有效关联下延续材料运动};
  \node[font=\trFontSmall,anchor=west] at (9.44,4.76)
    {新生表面从形成后记录};

  % ---- inter-column constraints ----
  % A single bidirectional connector per pair, placed in the empty band
  % between the column header and the first card.  The two labels describe
  % the two directions; no label covers a card border.
  \draw[trDoubleSmall] (4.53,9.10) -- (4.79,9.10);
  \draw[trDoubleSmall] (8.91,9.10) -- (9.17,9.10);

  % ---- common representation ----
  \path[fill=trBluePale,draw=trBlueDark,line width=0.85pt]
    (0.25,3.05) rectangle (13.55,4.32);
  \path[fill=trBlueDark] (0.25,3.84) rectangle (13.55,4.32);
  \node[font=\trFontHead,text=white] at (6.90,4.08)
    {时空一致的连续动态场景表征};
  \matrix[matrix of nodes,ampersand replacement=\&,
    nodes={inner sep=0pt,anchor=center,font=\trFontMath},
    column sep=0.8pt,row sep=0.6pt] at (6.90,3.50) {
    $S(t)=\bigl($ \& $\mathcal G(t)$ \& $,$ \& $\mathcal C(t)$ \& $,$ \&
      $\mathcal M(t)$ \& $,$ \& $\mathcal Z(t)$ \& $,$ \& $\mathcal E(t)$ \& $\bigr)$\\
     \& {\trFontSmall\color{trBlueDark} 几何} \& \& {\trFontSmall\color{trBlueDark} 外观} \& \&
      {\trFontSmall\color{trBlueDark} 连续运动} \& \& {\trFontSmall\color{trGreenDark} 空间归属} \& \&
      {\trFontSmall\color{trOchreDark} 拓扑结构} \& \\
  };
  % ================= action-conditioned prediction =================
  \draw[trDown] (5.40,3.02) -- (5.40,2.82);

  \path[fill=trVioletWash,draw=trVioletDark,line width=0.8pt]
    (2.30,1.12) rectangle (13.55,2.80);
  \path[fill=trVioletDark] (2.30,2.30) rectangle (13.55,2.80);
  \node[trNumOpen,draw=trVioletDark,text=trVioletDark,fill=white,
    minimum size=0.30cm] at (2.52,2.55) {4};
  \node[font=\trFontHead,text=white,anchor=west]
    at (2.72,2.55) {基于连续动态表征的动作后果预测};

  \ovActionCardOne{2.45}{4.96}{1.22}{2.18}{读取 $t_0$ 候选\\与动作对齐}
  \ovActionCardOne{5.26}{7.77}{1.22}{2.18}{分类型关系传播}
  \ovActionCardOne{8.07}{10.58}{1.22}{2.18}{逐候选连续推演}
  \ovActionCardOne{10.89}{13.40}{1.22}{2.18}{候选响应及范围}
  \draw[trFlow] (4.99,1.70) -- (5.23,1.70);
  \draw[trFlow] (7.80,1.70) -- (8.04,1.70);
  \draw[trFlow] (10.61,1.70) -- (10.85,1.70);

  % action and history inputs
  \path[fill=white,draw=trVioletMid,line width=0.6pt,rounded corners=1pt]
    (0.25,1.96) rectangle (2.12,2.52);
  \node[font=\trFontSubhead] at (1.18,2.24) {给定动作};
  \path[fill=white,draw=trVioletMid,line width=0.6pt,rounded corners=1pt]
    (0.25,1.32) rectangle (2.12,1.88);
  \node[font=\trFontSubhead] at (1.18,1.60) {已获交互历史};
  \draw[trFlow] (2.15,2.24) -- (2.42,2.24);
  \draw[trFlow] (2.15,1.60) -- (2.42,1.60);

  % offline evaluation, subsequent observations and representation update
  \ovCardOne{2.45}{5.45}{0.42}{0.86}{预测结果保存}
  \ovCardOne{5.85}{8.85}{0.42}{0.86}{后续观测对照}
  \ovCardOne{9.25}{12.25}{0.42}{0.86}{表征与历史更新}
  \draw[trFlow] (5.48,0.64) -- (5.82,0.64);
  \draw[trFlow] (8.88,0.64) -- (9.22,0.64);

  % prediction output feeds offline error evaluation through the gap below
  % the prediction frame.
  \draw[trFlow] (12.31,1.22) -- (12.31,1.00) -- (5.45,1.00) -- (5.45,0.87);
  \fill[trVioletDark] (12.31,1.22) circle (0.040);

  % new observation returns to the shared representation, bypassing initialisation.
  % The label is rotated so it runs alongside, not across, the dashed line.
  \draw[trFeedback] (12.28,0.38) -- (14.20,0.38) -- (14.20,3.12) -- (13.58,3.12);
  \node[rotate=90,font=\trFontTiny]
    at (14.42,2.07) {新增观测与表征修正};

  % ================= validation =================
  \path[fill=trBlueWash,draw=trBlueDark,line width=0.8pt]
    (0.25,-1.08) rectangle (13.55,0.12);
  \path[fill=trBlueDark] (0.25,-0.34) rectangle (13.55,0.12);
  \node[font=\trFontHead,text=white] at (6.90,-0.11)
    {公开数据与机制验证};
  \foreach \xa/\xb/\txt in {%
    0.35/4.55/{几何、运动与结构质量},%
    4.80/9.00/{响应预测与误差传递},%
    9.25/13.45/{可区分条件与适用范围}}{%
    \path[fill=white,draw=trBlueMid!55,line width=0.45pt]
      (\xa,-0.99) rectangle (\xb,-0.43);
    \node[font=\trFontSubhead]
      at ({(\xa+\xb)/2},-0.71) {\txt};
  }
\end{tikzpicture}
\endgroup
